\subsection{Hardware Setup and Implementation}

The experimental platform integrates multiple hardware components to enable real-time 
unified risk assessment and human-in-the-loop evaluation (Figure~\ref{fig:hardware_setup}).

\textbf{Computing Infrastructure.} The primary computation runs on a workstation 
equipped with an NVIDIA RTX-4090 GPU hosting the CARLA 0.9.14 simulation environment 
and Transfuser++ autonomous driving agent. Real-time processing operates at 10~Hz 
to balance computational efficiency with responsive risk assessment.

\textbf{Driver Monitoring System.} An NVIDIA AGX Orin edge computing device performs 
real-time driver state classification using a ResNet-50 convolutional neural network. 
The system processes camera input at 10~Hz and transmits driver state classifications 
(\textit{safe\_driving}, \textit{sleepy}, \textit{reaching\_back}, \textit{using\_phone}) 
with confidence scores via UDP packets to the main risk management system.

\textbf{Human Interface.} Manual vehicle control is provided through a Logitech G29 
force-feedback steering wheel with integrated pedals, mounted on a dedicated racing 
seat to simulate realistic driving ergonomics. This setup enables participants to 
perform natural driving behaviors while experiencing the risk management system's alerts.

\textbf{Audio-Visual Feedback.} The system delivers multi-modal alerts through 
high-quality speakers for audio warnings (beeps and text-to-speech) and a large 
display showing real-time risk visualization, ego vehicle speed, and traffic signal 
status. Alert audio is generated using system beep commands and \texttt{espeak} 
text-to-speech synthesis with optimized parameters for clarity and urgency.

\textbf{Network Architecture.} Communication between components utilizes standard 
Ethernet networking: the AGX Orin transmits driver state data via UDP to port 9999, 
while the main system broadcasts vehicle state and risk information for logging and 
analysis. This distributed architecture ensures minimal latency while maintaining 
system modularity for experimental control.

\begin{figure}[ht]
  \centering
  \includegraphics[width=0.9\columnwidth]{figs/hardware_setup.jpg}
  \caption{Experimental hardware setup showing the integrated workstation, 
           AGX Orin driver monitoring system, Logitech G29 steering wheel, 
           and audio-visual feedback components used for human-in-the-loop 
           evaluation of the unified risk management system.}
  \label{fig:hardware_setup}
\end{figure}

\textbf{System Integration.} The complete hardware stack enables end-to-end evaluation 
from perception through decision-making to human feedback. Synchronized data logging 
captures all system states, risk assessments, alert events, and participant responses 
at 10~Hz, providing comprehensive datasets for ablation analysis and performance 
validation. The modular design supports rapid reconfiguration for different 
experimental conditions via environment variables, facilitating systematic 
component ablation studies. 